\chapter{DEFINICIÓN DEL PROYECTO DE INVESTIGACIÓN}

\section{Definición del problema}

\subsection{Situación problemática}
En la enseñanza del Jiu-Jitsu Brasileño (BJJ) dentro de la academia \textit{Corpo e Mente}, los alumnos principiantes e intermedios enfrentan dificultades significativas para evaluar y perfeccionar su ejecución técnica de manera objetiva y continua. En el modelo tradicional, la enseñanza depende casi exclusivamente de la observación directa y presencial del profesor durante clases grupales numerosas. 

Esta dinámica genera los siguientes problemas críticos:
\begin{itemize}
    \item \textbf{Atención individualizada limitada:} En clases con alta afluencia de alumnos, el profesor no puede observar la totalidad de las repeticiones ejecutadas por cada alumno.
    \item \textbf{Retroalimentación diferida o ausente:} Si un error postural o de apalancamiento no es detectado en el instante preciso de la ejecución, el alumno tiende a automatizar patrones de movimiento incorrectos.
    \item \textbf{Evaluación subjetiva sin métricas cuantitativas:} La corrección tradicional carece de mediciones objetivas sobre ángulos articulares, alineación espinal o posición del cuerpo.
    \item \textbf{Riesgo incrementado de lesiones:} La repetición continua de técnicas con una biomecánica defectuosa expone a las articulaciones a sobrecargas y tensiones anómalas.
    \item \textbf{Progreso estancado fuera del aula:} El alumno no dispone de herramientas de autoevaluación guiadas fuera de los horarios de entrenamiento para comparar su desempeño con la técnica patrón enseñada por su profesor.
\end{itemize}

\subsection{Situación deseada}
La situación deseada consiste en implementar una solución tecnológica de apoyo para la academia \textit{Corpo e Mente} basada en una aplicación web que integre visión por computadora e inteligencia artificial generativa.

Con esta solución se logrará:
\begin{enumerate}
    \item Permitir al alumno grabar sus ejecuciones técnicas desde su dispositivo móvil en cualquier momento.
    \item Extraer automáticamente la postura corporal del alumno y compararla de forma analítica con la técnica patrón registrada por el profesor.
    \item Identificar desajustes articulares y generar reportes pedagógicos estructurados (análisis postural, riesgo de lesión y corrección paso a paso).
    \item Entregar retroalimentación objetiva, inmediata y fundamentada en literatura especializada de BJJ, acelerando la curva de aprendizaje técnico y reduciendo el riesgo de lesiones.
\end{enumerate}

\subsection{Objeto de investigación}
El objeto de investigación es el proceso de evaluación y corrección biomecánica postural de técnicas de Jiu-Jitsu Brasileño mediante visión por computadora y modelos de lenguaje orquestados bajo arquitecturas de recuperación de información.

\subsection{Alcance}

\subsubsection{Alcance Funcional}
\begin{itemize}
    \item \textbf{Gestión de Usuarios y Roles:} Perfiles para Alumnos (evaluación técnica) y Profesores (administración de técnicas patrón y fuentes de conocimiento). El sistema implementa autenticación con PBKDF2-HMAC-SHA256 y validación de roles contra la tabla usuarios (rol IN ('alumno', 'profesor')).
    \item \textbf{Registro de Técnicas Patrón:} Carga de videos de referencia del profesor y extracción de posturas corporales patrón.
    \item \textbf{Evaluación Biomecánica:} Comparación postural entre el video del alumno y el patrón del profesor para calcular desajustes angulares.
    \item \textbf{Búsqueda en Literatura Técnica:} Procesamiento e ingesta semántica de manuales y libros de BJJ para fundamentar la retroalimentación.
    \item \textbf{Síntesis Pedagógica Estructurada:} Generación de reportes explicativos organizados en Análisis Postural, Riesgo de Lesión, Guía Paso a Paso y Resumen Ejecutivo.
    \item \textbf{Interfaz Web de Usuario:} Aplicación accesible para la carga de videos, consulta de resultados e historial de avance.
\end{itemize}

\subsubsection{Alcance Tecnológico}
\begin{itemize}
    \item \textbf{Aplicación Web y Servidor de Procesamiento:} Arquitectura cliente-servidor orientada a servicios para la captura de video, procesamiento de imágenes y almacenamiento de información.
    \item \textbf{Modelos de Inteligencia Artificial:} YOLO26x-Pose + YOLO26x-Depth para estimación de pose 3D (17 keypoints COCO), Qwen3-VL-Embedding-2B para vectorización densa de 2048 dimensiones, Gemini 2.5 Flash para síntesis pedagógica estructurada en JSON.
\end{itemize}

\subsubsection{Exclusiones}
\begin{itemize}
    \item El sistema no sustituye la evaluación presencial formal ni el criterio del profesor para la graduación de cinturones.
    \item No realiza diagnósticos médicos o clínicos fisioterapéuticos de rehabilitación.
\end{itemize}

\subsection{Justificación}

\subsubsection{Justificación Técnica}
El proyecto demuestra la aplicación convergente de tecnología moderna en ingeniería de software, visión artificial e inteligencia artificial generativa, organizada bajo patrones de diseño de software y principios de arquitectura limpia para asegurar mantenibilidad y calidad.

\subsubsection{Justificación Práctica y Social}
Brinda a los estudiantes de la academia \textit{Corpo e Mente} un asistente técnico disponible las 24 horas que democratiza el aprendizaje guiado, reduce la frustración en etapas iniciales y promueve una práctica deportiva segura al corregir malas posturas.

\subsubsection{Justificación Institucional}
Posiciona a la academia \textit{Corpo e Mente} a la vanguardia tecnológica en la enseñanza de artes marciales en la región, incorporando análisis biomecánico asistido por inteligencia artificial en su propuesta pedagógica.

\newpage

\section{Objetivos}

\subsection{Objetivo General}
Desarrollar un sistema de apoyo para el análisis biomecánico de movimientos en Jiu-Jitsu Brasileño para la academia \textit{Corpo e Mente}, mediante visión por computadora e inteligencia artificial, que permita evaluar objetivamente las ejecuciones técnicas de los alumnos y proporcionar retroalimentación pedagógica estructurada en tiempo real.

\subsection{Objetivos Específicos}
\begin{enumerate}
    \item Analizar y especificar los requerimientos del sistema y casos de uso para el registro de patrones técnicos y la evaluación biomecánica.
    \item Implementar el módulo de visión por computadora para la detección y extracción de posturas corporales a partir de videos de entrenamiento.
    \item Desarrollar el algoritmo de comparación biomecánica para computar las diferencias angulares articulares entre la ejecución del estudiante y la técnica patrón del instructor.
    \item Construir el motor de búsqueda semántica para recuperar información explicativa desde literatura especializada de Jiu-Jitsu Brasileño.
    \item Integrar el modelo de lenguaje generativo para elaborar reportes de retroalimentación pedagógica en análisis postural, riesgo de lesión, guía paso a paso y resumen ejecutivo.
    \item Desarrollar una interfaz web intuitiva y responsiva que permita la interacción eficiente entre profesores y estudiantes.
    \item Validar la solución mediante pruebas unitarias, de integración y funcionales.
\end{enumerate}

\newpage

\section{Metodología}

\subsection{Ingeniería de Software (Proceso Unificado)}
Para el desarrollo del software se adopta la metodología del Proceso Unificado (UP), guiada por casos de uso, centrada en la arquitectura de software e iterativa e incremental. Su ciclo de vida comprende cuatro fases:

\begin{enumerate}
    \item \textbf{Fase de Inicio (Inception):} Definición de los objetivos del proyecto, delimitación del alcance y evaluación de factibilidad.
    \item \textbf{Fase de Elaboración (Elaboration):} Especificación detallada de casos de uso, diseño de la arquitectura en capas y prototipado del motor biomecánico.
    \item \textbf{Fase de Construcción (Construction):} Desarrollo iterativo de los componentes del servidor (FastAPI), modelos de inteligencia artificial (YOLO26x en Colab GPU, Qwen3-VL-Embedding-2B, Gemini 2.5 Flash), pipeline RAG (Qdrant + PostgreSQL BCNF) e interfaz web progresiva (PWA).
    \item \textbf{Fase de Transición (Transition):} Pruebas del sistema, optimización de componentes y despliegue final para la academia \textit{Corpo e Mente}.
\end{enumerate}

\subsection{Gestión del Proyecto (Scrum)}
La gestión operativa del proyecto se alinea con el marco ágil \textbf{Scrum}, organizando el trabajo en iteraciones cortas denominadas Sprints:

\begin{itemize}
    \item \textbf{Planificación Temporal y Cronograma:} El desarrollo del proyecto contempla una duración total de un (1) año, iniciando las actividades de investigación y desarrollo en febrero de 2026 y finalizando con el despliegue y validación en diciembre de 2026.
    \item \textbf{Product Backlog:} Lista priorizada de requisitos y funcionalidades del sistema.
    \item \textbf{Sprints:} Ciclos de trabajo de 2 semanas para la entrega de incrementos funcionales probados del software.
    \item \textbf{Aseguramiento de Calidad:} Desarrollo guiado por pruebas para validar el correcto funcionamiento de la lógica de negocio e interfaces.
    \item \textbf{Control de Versiones:} Gestión del código fuente mediante control de versiones y gestión de entornos de trabajo.
\end{itemize}
